% GENERATED FILE. Edit canonical lessons, then run npm run generate:books.
% canonical-source-hash: fnv1a64:2272a4edea7c5341
% canonical-lessons: SA-C397-jvalayati, SA-C397-dipasikha, SA-C397-karpurah, SA-C397-stotram, SA-C397-stauti

\chapter{Lamps and Hymns}
\label{ch:sa-lamps-and-hymns}
\hlchaptermodality{397}

\begin{chapteropening}
\textbf{By the end of this chapter:} \emph{I can say what is lit and sung at evening worship.}

Complete the last lesson of chapter 397: I can say what is lit and sung at evening worship.
\end{chapteropening}

\section[\texorpdfstring{\sk{ज्वालयति} (jvālayati)}{ज्वालयति (jvālayati)}]{\sk{ज्वालयति} (jvālayati) — lights; kindles}

\label{lesson:SA-C397-jvalayati}



\begin{quote}
Before the new one: say the Sanskrit for an inner sanctum, then the Sanskrit for a devotee.
\end{quote}

\subsection*{You'll want to know: \sk{ज्वालयति}}
\textbf{\sk{ज्वालयति}} — \emph{jvālayati} — \textquotedblleft{}lights; kindles\textquotedblright{}.

\begin{quote}
\sk{माता} \sk{सायं} \sk{दीपं} \sk{ज्वालयति।}
\end{quote}

\begin{quote}
\emph{mātā sāyaṃ dīpaṃ jvālayati.}
\end{quote}

\begin{quote}
Mother lights a lamp in the evening.
\end{quote}

Past (she …): \sk{सा} \sk{दीपं} \sk{ज्वालितवती}, \emph{sā dīpaṃ jvālitavatī}.

\begin{cousinweb}
\textbf{\sk{ज्वालयति}} (\emph{jvālayati}) is built on the root √jval, \textquotedblleft{}to blaze\textquotedblright{}, which gives \emph{jvalati}, \textquotedblleft{}burns\textquotedblright{}; MW gives its causative, \textquotedblleft{}makes blaze\textquotedblright{}, as to set on fire, light or kindle, and KRM gives the causative participle \emph{jvālita} (MW).
\end{cousinweb}

\begin{grammarlens}[title={What you've built}]
One more word for getting around.
\end{grammarlens}

\subsection*{Your turn}
\begin{itemize}
\raggedright
  \item \emph{Say it:} \emph{jvālayati}
  \item \emph{Say it:} \emph{jvālayati}, once more
  \item \emph{Say it:} \emph{mātā sāyaṃ dīpaṃ jvālayati.} — Mother lights a lamp in the evening.
  \item \emph{Recall:} say the Sanskrit for an inner sanctum, then the Sanskrit for a devotee, then say \emph{jvālayati} again
\end{itemize}

\begin{tcolorbox}[breakable,colback=teal!4,colframe=teal!35!black,title={Before you move on}]
What does \sk{ज्वालयति} mean? (\textquotedblleft{}Lights; kindles\textquotedblright{}.) Say it once more.
\end{tcolorbox}
\section[\texorpdfstring{\sk{दीपशिखा} (dīpaśikhā)}{दीपशिखा (dīpaśikhā)}]{\sk{दीपशिखा} (dīpaśikhā) — a lamp's flame}

\label{lesson:SA-C397-dipasikha}



\begin{quote}
Before the new one: say the Sanskrit for a devotee, then the Sanskrit for kindles.
\end{quote}

\subsection*{You'll want to know: \sk{दीपशिखा}}
\textbf{\sk{दीपशिखा}} — \emph{dīpaśikhā} — \textquotedblleft{}a lamp's flame\textquotedblright{}.

\begin{quote}
\sk{वायुना} \sk{दीपशिखा} \sk{चलति।}
\end{quote}

\begin{quote}
\emph{vāyunā dīpaśikhā calati.}
\end{quote}

\begin{quote}
The lamp's flame flickers in the wind.
\end{quote}

A feminine noun.

\begin{cousinweb}
\textbf{\sk{दीपशिखा}} (\emph{dīpaśikhā}) is built from \emph{dīpa}, \textquotedblleft{}a lamp\textquotedblright{}, and \emph{śikhā}, \textquotedblleft{}a tuft, a point of flame\textquotedblright{}; MW gives the compound as the flame of a lamp (MW).
\end{cousinweb}

\begin{grammarlens}[title={What you've built}]
One more word for getting around.
\end{grammarlens}

\subsection*{Your turn}
\begin{itemize}
\raggedright
  \item \emph{Say it:} \emph{dīpaśikhā}
  \item \emph{Say it:} \emph{dīpaśikhā}, once more
  \item \emph{Say it:} \emph{vāyunā dīpaśikhā calati.} — The lamp's flame flickers in the wind.
  \item \emph{Recall:} say the Sanskrit for a devotee, then the Sanskrit for kindles, then say \emph{dīpaśikhā} again
\end{itemize}

\begin{tcolorbox}[breakable,colback=teal!4,colframe=teal!35!black,title={Before you move on}]
What does \sk{दीपशिखा} mean? (\textquotedblleft{}A lamp's flame\textquotedblright{}.) Say it once more.
\end{tcolorbox}
\section[\texorpdfstring{\sk{कर्पूरः} (karpūraḥ)}{कर्पूरः (karpūraḥ)}]{\sk{कर्पूरः} (karpūraḥ) — camphor}

\label{lesson:SA-C397-karpurah}



\begin{quote}
Before the new one: say the Sanskrit for kindles, then the Sanskrit for a lamp's flame.
\end{quote}

\subsection*{You'll want to know: \sk{कर्पूरः}}
\textbf{\sk{कर्पूरः}} — \emph{karpūraḥ} — \textquotedblleft{}camphor\textquotedblright{}.

\begin{quote}
\sk{दीपे} \sk{कर्पूरः} \sk{ज्वलति।}
\end{quote}

\begin{quote}
\emph{dīpe karpūraḥ jvalati.}
\end{quote}

\begin{quote}
Camphor burns in the lamp.
\end{quote}

A masculine noun.

\begin{cousinweb}
\textbf{\sk{कर्पूरः}} (\emph{karpūraḥ}) is the word MW gives for camphor; MW records only a commentator's guess at a root √kṛp, so its origin is uncertain (MW).
\end{cousinweb}

\begin{grammarlens}[title={What you've built}]
One more word for getting around.
\end{grammarlens}

\subsection*{Your turn}
\begin{itemize}
\raggedright
  \item \emph{Say it:} \emph{karpūraḥ}
  \item \emph{Say it:} \emph{karpūraḥ}, once more
  \item \emph{Say it:} \emph{dīpe karpūraḥ jvalati.} — Camphor burns in the lamp.
  \item \emph{Recall:} say the Sanskrit for kindles, then the Sanskrit for a lamp's flame, then say \emph{karpūraḥ} again
\end{itemize}

\begin{tcolorbox}[breakable,colback=teal!4,colframe=teal!35!black,title={Before you move on}]
What does \sk{कर्पूरः} mean? (\textquotedblleft{}Camphor\textquotedblright{}.) Say it once more.
\end{tcolorbox}
\section[\texorpdfstring{\sk{स्तोत्रम्} (stotram)}{स्तोत्रम् (stotram)}]{\sk{स्तोत्रम्} (stotram) — a hymn of praise}

\label{lesson:SA-C397-stotram}



\begin{quote}
Before the new one: say the Sanskrit for a lamp's flame, then the Sanskrit for camphor.
\end{quote}

\subsection*{You'll want to know: \sk{स्तोत्रम्}}
\textbf{\sk{स्तोत्रम्}} — \emph{stotram} — \textquotedblleft{}a hymn of praise\textquotedblright{}.

\begin{quote}
\sk{बालाः} \sk{सायं} \sk{स्तोत्रं} \sk{गायन्ति।}
\end{quote}

\begin{quote}
\emph{bālāḥ sāyaṃ stotraṃ gāyanti.}
\end{quote}

\begin{quote}
The children sing a hymn in the evening.
\end{quote}

A neuter noun.

\begin{cousinweb}
\textbf{\sk{स्तोत्रम्}} (\emph{stotram}) is built on the root √stu, \textquotedblleft{}to praise\textquotedblright{}, with the ending \emph{-tra}; MW gives \emph{stotra} as praise, a hymn of praise (MW).
\end{cousinweb}

\begin{grammarlens}[title={What you've built}]
One more word for getting around.
\end{grammarlens}

\subsection*{Your turn}
\begin{itemize}
\raggedright
  \item \emph{Say it:} \emph{stotram}
  \item \emph{Say it:} \emph{stotram}, once more
  \item \emph{Say it:} \emph{bālāḥ sāyaṃ stotraṃ gāyanti.} — The children sing a hymn in the evening.
  \item \emph{Recall:} say the Sanskrit for a lamp's flame, then the Sanskrit for camphor, then say \emph{stotram} again
\end{itemize}

\begin{tcolorbox}[breakable,colback=teal!4,colframe=teal!35!black,title={Before you move on}]
What does \sk{स्तोत्रम्} mean? (\textquotedblleft{}A hymn of praise\textquotedblright{}.) Say it once more.
\end{tcolorbox}
\section[\texorpdfstring{\sk{स्तौति} (stauti)}{स्तौति (stauti)}]{\sk{स्तौति} (stauti) — sings the praises of}

\label{lesson:SA-C397-stauti}



\begin{quote}
Before the new one: say the Sanskrit for camphor, then the Sanskrit for a hymn of praise.
\end{quote}

\subsection*{You'll want to know: \sk{स्तौति}}
\textbf{\sk{स्तौति}} — \emph{stauti} — \textquotedblleft{}sings the praises of\textquotedblright{}.

\begin{quote}
\sk{शिष्यः} \sk{श्लोकेन} \sk{गुरुं} \sk{स्तौति।}
\end{quote}

\begin{quote}
\emph{śiṣyaḥ ślokena guruṃ stauti.}
\end{quote}

\begin{quote}
The disciple praises his teacher in a verse.
\end{quote}

Past (he …): \sk{सः} \sk{गुरुं} \sk{स्तुतवान्}, \emph{saḥ guruṃ stutavān}.

\begin{cousinweb}
\textbf{\sk{स्तौति}} (\emph{stauti}) comes from the root √stu, \textquotedblleft{}to praise, to hymn\textquotedblright{}, which also gives \emph{stotra}, \textquotedblleft{}a hymn\textquotedblright{}; MW gives the participle \emph{stuta}, \textquotedblleft{}praised, hymned\textquotedblright{} (MW).
\end{cousinweb}

\begin{grammarlens}[title={What you've built}]
One more word for getting around.
\end{grammarlens}

\subsection*{Your turn}
\begin{itemize}
\raggedright
  \item \emph{Say it:} \emph{stauti}
  \item \emph{Say it:} \emph{stauti}, once more
  \item \emph{Say it:} \emph{śiṣyaḥ ślokena guruṃ stauti.} — The disciple praises his teacher in a verse.
  \item \emph{Recall:} say the Sanskrit for camphor, then the Sanskrit for a hymn of praise, then say \emph{stauti} again
\end{itemize}

\begin{tcolorbox}[breakable,colback=teal!4,colframe=teal!35!black,title={Before you move on}]
What does \sk{स्तौति} mean? (\textquotedblleft{}Sings the praises of\textquotedblright{}.) Say it once more.
\end{tcolorbox}
